\section{Исследование производительности и качества}

\subsection{Дневник выполнения работы}

В соответствии с заданием, ниже представлен протокол выполнения работы по дням.

\begin{center}
\begin{tabular}{|p{3cm}|p{7cm}|p{4cm}|}
\hline
\textbf{Дата} & \textbf{Описание проделанной работы} & \textbf{Результат} \\ \hline
10.04.2026 & Реализация первой версии словаря на основе Patricia trie. & Корректная работа операций вставки, поиска и удаления. \\ \hline
11.04.2026 & Подготовка набора тестов для проверки производительности на разных объёмах входных данных. & Получены замеры для 100\,000, 200\,000 и 400\,000 операций. \\ \hline
12.04.2026 & Анализ программы с помощью \texttt{gprof} и \texttt{Valgrind}. & Найдены лишние копирования строк и повышенная нагрузка на менеджер памяти. \\ \hline
13.04.2026 & Оптимизация: использование семантики перемещения и доработка очистки дерева. & Получена версия с меньшим числом аллокаций и более стабильным временем работы. \\ \hline
14.04.2026 & Повторный запуск бенчмарков на тех же размерах входа. & Подтверждён линейный рост времени работы относительно числа операций. \\ \hline
\end{tabular}
\end{center}

\subsection{Методика бенчмарка}

Для проверки масштабируемости были проведены замеры на трёх размерах входных данных:
100\,000, 200\,000 и 400\,000 операций. В каждом тесте выполнялась одинаковая последовательность действий:
вставка ключей, поиск всех добавленных ключей и последующая очистка структуры данных.

Такой набор размеров выбран специально: при переходе от $N$ к $2N$ количество операций удваивается.
Если алгоритм работает линейно, то общее время выполнения также должно увеличиваться примерно в два раза,
а среднее время одной операции должно оставаться почти постоянным.

\subsection{Анализ времени работы}

В таблице приведены результаты замеров времени для двух версий программы. Значение
$T(2N) / T(N)$ показывает, насколько увеличилось время работы при удвоении размера входа.

\begin{center}
\begin{tabular}{|r|r|r|r|r|}
\hline
\textbf{Число операций} &
\textbf{v1, сек.} &
\textbf{v1, мкс/операция} &
\textbf{v2, сек.} &
\textbf{v2, мкс/операция} \\ \hline
100\,000 & 0{,}124 & 1{,}24 & 0{,}108 & 1{,}08 \\ \hline
200\,000 & 0{,}249 & 1{,}25 & 0{,}216 & 1{,}08 \\ \hline
400\,000 & 0{,}501 & 1{,}25 & 0{,}434 & 1{,}09 \\ \hline
\end{tabular}
\end{center}

\begin{center}
\begin{tabular}{|r|r|r|}
\hline
\textbf{Переход} & \textbf{v1: $T(2N)/T(N)$} & \textbf{v2: $T(2N)/T(N)$} \\ \hline
100\,000 $\rightarrow$ 200\,000 & 2{,}01 & 2{,}00 \\ \hline
200\,000 $\rightarrow$ 400\,000 & 2{,}01 & 2{,}01 \\ \hline
\end{tabular}
\end{center}

Полученные значения близки к 2, то есть при удвоении количества операций время выполнения также
увеличивается примерно в два раза. При этом среднее время одной операции остаётся почти неизменным:
около 1{,}24--1{,}25 мкс для первой версии и около 1{,}08--1{,}09 мкс для второй версии.
Это подтверждает линейную зависимость времени работы от числа операций.

\subsection{Анализ профилирования времени (\texttt{gprof})}

Профилирование с помощью \texttt{gprof} показало, что в первой версии заметная часть времени
тратилась на вспомогательные операции со строками. После оптимизации нагрузка сместилась на основные
операции структуры данных: вставку, поиск и получение бита ключа.

\textbf{Версия 1 до оптимизации:}
\begin{verbatim}
  %   cumulative   self              self     total
 time   seconds   seconds    calls  ms/call  ms/call  name
 27.78      0.03     0.03   100000     0.00     0.00  Patricia::toLower(...)
 11.11      0.07     0.01 25867084     0.00     0.00  __gnu_cxx::__normal_iterator...
 11.11      0.08     0.01 12933542     0.00     0.00  bool __gnu_cxx::operator!=...
\end{verbatim}

\textbf{Версия 2 после оптимизации:}
\begin{verbatim}
  %   cumulative   self              self     total
 time   seconds   seconds    calls  ms/call  ms/call  name
 25.00      0.03     0.03   100000     0.00     0.00  Patricia::toLower(...)
 16.67      0.07     0.02    69848     0.00     0.00  Patricia::search(...) const
  8.33      0.08     0.01  4365504     0.00     0.00  Patricia::getBit(...)
\end{verbatim}

\subsection{Анализ потребления памяти (\texttt{Valgrind})}

В ходе тестирования с помощью \texttt{Valgrind} были получены следующие данные:

\begin{itemize}
    \item \textbf{Версия 1}: всего аллокаций: 400\,375. Общий объём: 52\,955\,238 байт.
    \item \textbf{Версия 2}: всего аллокаций: 362\,574. Общий объём: 48\,805\,174 байт.
\end{itemize}

\textbf{Вывод}: оптимизация позволила сократить количество выделений памяти и уменьшить нагрузку
на менеджер памяти при массовых операциях вставки и поиска. При этом серия замеров на
100\,000, 200\,000 и 400\,000 операций показывает, что время работы растёт линейно относительно
размера входных данных.

\subsection{Выводы о найденных недочётах}

\begin{enumerate}
    \item \textbf{Копирование строк}: передача ключа по значению приводила к появлению лишних временных
    объектов. Исправлено с помощью \texttt{std::move}.
    \item \textbf{Рекурсивный \texttt{clear}}: при большой глубине дерева рекурсивная очистка могла
    приводить к переполнению стека. Для устранения проблемы был реализован итеративный метод очистки
    с использованием очереди.
    \item \textbf{Недостаточный набор размеров в бенчмарке}: один запуск на 100\,000 операций не позволял
    сделать вывод о характере роста времени. После добавления тестов на 200\,000 и 400\,000 операций
    стало видно, что зависимость времени от числа операций близка к линейной.
\end{enumerate}

\pagebreak
